\documentclass[10pt,conference,compsocconf]{IEEEtran}

\usepackage{hyperref}
\usepackage{graphicx}
\usepackage{amsmath}
\usepackage{booktabs}

\begin{document}
\title{Predicting Coronary Heart Disease from Lifestyle Survey Data\\with Regularized Logistic Regression}

\author{
  Team members (fill in)\\
  \textit{CS-433 Machine Learning, EPFL, Switzerland}
}

\maketitle

\begin{abstract}
We predict whether a respondent of the 2015 Behavioral Risk Factor
Surveillance System (BRFSS) has been diagnosed with coronary heart disease or
myocardial infarction (MICHD) from 321 survey answers. Using only NumPy, we
compare linear and logistic regression baselines and build a feature pipeline
that one-hot encodes categorical answers, keeps missingness as information,
cleans survey answer codes and discretizes numeric variables. Combined with
$L_2$-regularized logistic regression solved by Newton's method and an
F1-optimal decision threshold, our model reaches an F1 score of
$0.440\pm0.008$ in 5-fold cross-validation. An ablation study shows that
choosing the decision threshold matters far more than any feature
engineering step: it raises F1 from 0.27 to 0.44.
\end{abstract}

\section{Introduction}
Cardiovascular diseases are a leading cause of death worldwide, and
identifying at-risk individuals from cheap survey data could support
prevention. The task is binary classification: given a respondent's answers,
predict MICHD ($y=1$) or not ($y=-1$). The data contain 328{,}135 training and
109{,}379 test respondents. Only 8.8\% of training respondents are positive,
and the competition metric is the F1 score of the positive class. Both facts
shape our design choices.

\section{Data and Feature Engineering}
\label{sec:features}
\textbf{Exploratory analysis.} Of the 321 columns, 239 contain missing values
and 147 are more than half missing. This is mostly by design: BRFSS skips
follow-up questions when they do not apply (e.g.\ diabetes questions for
non-diabetics). 261 columns have at most 60 distinct values and encode
categorical answers. Raw questions use special codes (e.g.\ 77 = ``don't
know'', 99 = ``refused'', 88 = ``none'') that would distort a numeric
treatment. The features most correlated with the label are general health,
age, employment status (retired), high blood pressure and diabetes.

\textbf{Pipeline.} All statistics below are fitted on training data only and
re-fitted inside every cross-validation fold.
\begin{enumerate}
\item \emph{Column filtering:} we drop 18 survey bookkeeping columns
(interview date, sequence numbers, sampling weights, phone screening) as well
as constant columns.
\item \emph{Categorical columns} ($\le 30$ distinct values) are one-hot encoded,
keeping every category that appears in $\ge 0.1\%$ of training rows. We add
an indicator for ``missing'' because a skipped question is informative.
\item \emph{Numeric columns} ($>30$ values): in raw questions, ``none'' codes
become 0 and ``don't know/refused'' codes become missing. A code is cleaned
only if it lies above every regular answer of that question, so that
e.g.\ 7 days of exercise is not removed. Missing values are imputed with the
training median and flagged with an indicator. Values are clipped to the
$[0.5, 99.5]$ percentiles and standardized. We also add indicators for 10
quantile bins, which lets the linear model learn a piecewise-constant,
non-monotone effect of e.g.\ age or BMI.
\end{enumerate}
The final design matrix has 1{,}733 features plus a bias column.

\section{Models and Training}
\textbf{Basic methods.} We implemented the six required methods: least squares
via the normal equations, ridge regression, MSE gradient descent and SGD, and
(regularized) logistic regression with gradient descent. We use ridge
regression (on targets in $\{-1,1\}$) and gradient-descent logistic
regression as baselines.

\textbf{Final model.} We minimize the regularized logistic loss
\[
\mathcal{L}(w)=\frac{1}{N}\sum_{i=1}^{N}\left[\log\!\left(1+e^{x_i^\top w}\right)-y_i x_i^\top w\right]+\lambda\lVert w\rVert^2
\]
with $y_i\in\{0,1\}$, leaving the bias unpenalized. Plain gradient descent
converges slowly on 1{,}700 correlated indicator features. We therefore use
Newton's method, $w\leftarrow w-t\,H^{-1}\nabla\mathcal{L}$ with
$H=\frac1N X^\top\mathrm{diag}(p_i(1-p_i))X+2\lambda I$, and a backtracking
line search on $t$. It converges in about 8 iterations ($\sim$10\,s on the
full data). Regularization is necessary: with $\lambda=0$ the Hessian is
singular because the one-hot blocks are collinear.

\textbf{Decision threshold.} With 8.8\% positives, the usual cut-off
$p>0.5$ predicts very few positives and gives poor F1. We sort the training
scores and pick the threshold that maximizes training F1 exactly; it is
$p\approx0.21$.

\textbf{Evaluation protocol.} We use 5-fold cross-validation with the same
random folds for every configuration. In each fold, the feature pipeline, the
model and the threshold are all fitted on the four training folds and
evaluated once on the held-out fold. We report the mean (and standard
deviation) of F1 over folds, together with accuracy and ROC AUC.

\section{Results}
Table~\ref{tab:ablation} starts from baselines and adds one modification at a
time. \emph{Raw} features are median-imputed, code-cleaned and standardized
original columns.

\begin{table}[htbp]
  \centering
  \caption{5-fold CV results. All rows except the last use the tuned threshold.}
  \label{tab:ablation}
  \begin{tabular}{@{}lccc@{}}
    \toprule
    Configuration & F1 & Acc. & AUC \\
    \midrule
    Ridge regression, raw & 0.429 & 0.875 & 0.858 \\
    Logistic, GD (300 it., $\gamma{=}0.5$), raw & 0.416 & 0.864 & 0.850 \\
    Logistic, Newton, raw & 0.436 & 0.877 & 0.864 \\
    \quad + drop admin columns & 0.436 & 0.880 & 0.864 \\
    \quad + one-hot categoricals & 0.439 & 0.876 & 0.866 \\
    \quad + quantile bins (\textbf{final}) & \textbf{0.440} & 0.880 & 0.865 \\
    Ridge regression, final features & 0.432 & 0.876 & 0.858 \\
    \midrule
    Final model, threshold $p=0.5$ & 0.267 & -- & -- \\
    \bottomrule
  \end{tabular}
\end{table}

\textbf{Threshold.} Moving the threshold from 0.5 to the F1-optimal value is
by far the most important choice ($+0.17$ F1). On a held-out 20\% split,
recall rises from 0.17 to 0.55 while precision falls from 0.60 to 0.37.
Accuracy drops from 0.916 to 0.875, which shows why accuracy alone is
misleading under class imbalance.

\textbf{Optimizer.} With the same objective, 300 steps of gradient descent
reach only F1 0.416, compared with 0.436 for Newton's method. The gap comes
from incomplete optimization, not from the model.

\textbf{Features.} One-hot encoding and binning add $+0.004$ F1 and $+0.001$
AUC over the raw logistic model. This is consistent across folds but
comparable to the fold-to-fold standard deviation ($\approx0.006$). The raw
features already capture most of the signal because many BRFSS variables are
ordinal with a roughly monotone risk. Dropping bookkeeping columns barely
changes performance, but it makes the model rely on health information
rather than on survey logistics.

\textbf{Logistic vs.\ least squares.} With identical features, logistic
regression beats ridge regression by $0.008$ F1 and $0.007$ AUC. Squared loss
penalizes confidently correct predictions and is therefore a poor surrogate
for classification under class imbalance.

\textbf{Regularization.} Sweeping $\lambda\in\{10^{-6},\dots,10^{-2}\}$, F1
is flat (0.439--0.440) for $\lambda\in[10^{-5},10^{-4}]$. It decreases for
stronger penalties (0.436 at $10^{-3}$, 0.422 at $10^{-2}$). We use
$\lambda=10^{-4}$, which has the best AUC. With 1{,}700 features and over
260k training rows, the model is far from overfitting: training F1 (0.449)
is close to validation F1.

\textbf{Test set.} Trained on all data, the model predicts 12.7\% of test
respondents as positive. Its leaderboard score is: F1 = \emph{(fill in)},
accuracy = \emph{(fill in)}.

\section{Discussion}
A well-optimized linear model with careful handling of survey codes and
missingness gives a strong baseline. Most of the achievable F1 comes from
aligning the decision rule with the metric rather than from model
complexity. The ranking quality (AUC $\approx0.865$) changes little across
our feature sets, which suggests the remaining error is largely irreducible
given these self-reported answers. Natural next steps are explicit
interactions (e.g.\ age $\times$ sex) and non-linear models; we leave them to
future work.

\end{document}
